\section{The family and its geometric parameter}\label{sec:family}
\subsection{Conventions}
\begin{definition}[FAM conventions]\label{def:conventions}
We use $B=\Pp^1_t$, $p_1=0$, $p_2=1$, $p_0=\infty$, and
$B^\circ=B\setminus\{0,1,\infty\}$. Base meridians are counterclockwise
in a local holomorphic coordinate. The analytic generators are clockwise
root deck transformations, $s_j(g_jz)=e^{-2\pi i/m_j}s_j(z)$, with
$g_1g_2g_0=1$. A positive base loop ends at the inverse deck translate;
identifying its endpoint with the initial fibre applies $g_j$, giving
homology monodromy $A_j$. The retained multiplicities are $m_1=3,m_2=4$;
$f^*p_j=m_jS_j$ includes the nonreduced structure, whereas $S_j$ is reduced.
The cusp fibre $W=f^{-1}(\infty)$ is reduced.
\[
 \Lambda=\Z\langle\gh,\uh,\wh,\dhat\rangle,\qquad
 \Pi_c=\begin{pmatrix}6\mu&\tau&1&0\\\beta^{\partt}+c&\mu&0&1\end{pmatrix},
 \qquad w=e^{-\pi ic}.
\]
We fix $(\beta^{\partt}+\tau)(p_0)=0$. Write
$D_{\partt}=\Im\beta^{\partt}-6(\Im\mu)^2/\Im\tau$ and
$M=\sup D_{\partt}$. Engel's line bundle is
$M_S=\OO_S(P-O)$; his indices $0,1,\infty$ correspond to our $1,2,0$.
In his common integral basis, the columns of the following matrix are
our ordered basis vectors:
\[
 P_{\rm mark}=\begin{pmatrix}
 0&-2&-1&-2\\2&-1&0&-1\\0&2&1&3\\-1&0&0&0
 \end{pmatrix},\quad
 T_0P_{\rm mark}=P_{\rm mark}A_1,\quad
 T_1P_{\rm mark}=P_{\rm mark}A_2,\quad
 \delta_{\rm E}=P_{\rm mark}\dhat.
\]
For our finite twist vectors $v_1,v_2$, his positive-root affine vectors
are $a_0=-P_{\rm mark}v_1/3$, $a_1=-P_{\rm mark}v_2/4$.
In additive fibre coordinates $(\xi,\eta)$, put $v=e^{2\pi i\eta}$.
The first period acts by
$(\xi,v)\mapsto(\xi+6\mu,e^{2\pi i\beta_c}v)$; Engel's contracting
generator is its inverse, over $-6P$ with multiplier $e^{-2\pi i\beta_c}$.
We choose his reference linearization to be this transported inverse
period lift at $c=0$. Then $\lambda_0=1$ and $u=w^2$.
The comparison maps and the effect of changing this reference are
proved below. Homology has integral coefficients unless stated
otherwise; tori and the threefold have their complex orientations.

\end{definition}
\uses{The period construction of \cite[Sections 2--6]{AC26}; \cite[Sections 2--5]{E26}.}

Throughout this section $\mathcal P=\{c:\operatorname{Im}c<-M\}$, and $X_c$ means the actual glued space with $v_1=\varepsilon=(1,2,-4,0)^t$, $v_2=-\varepsilon'=(-1,-3,3,0)^t$ and the untwisted cusp collar. We write $d=\widehat\delta$, $e=(0,1)^t$, $E=d\gamma$, and use column vectors. Our peripheral generators are clockwise: $s_j(g_jz)=e^{-2\pi i/m_j}s_j(z)$ and $g_1g_2g_0=1$. Thus $s(g_0z)=s(z)-1$ at the cusp. All period comparisons use
\[
 A_1=\begin{pmatrix}1&0&0&0\\6&0&1&0\\-6&-1&-1&0\\-2&1&0&1\end{pmatrix},\quad
 A_2=\begin{pmatrix}1&0&0&0\\0&0&-1&0\\-6&1&0&0\\3&0&1&1\end{pmatrix},\qquad
 R_g(z)\Pi_c(z)=\Pi_c(gz)A_g.
\]
Here $R_1=\left(\begin{smallmatrix}-1/\tau&0\\(1-\mu)/\tau&1\end{smallmatrix}\right)$, $R_2=\left(\begin{smallmatrix}1/\tau&0\\-\mu/\tau&1\end{smallmatrix}\right)$, and $R_0=I$. The distinction between source and target markings in each equation below is retained by its displayed period equation.


The construction, the smooth manifold $S^6$, its algebraic reduction, its canonical bundle and the connected automorphism group are due to \cite[Main Theorem; Sections 3--9]{AC26}. The rational elliptic surface, its line bundle and inverse-translation quotient are those of \cite[Section 2]{E26}. The equality $u=w^2$ uses the transported reference in Definition~\ref{def:conventions}; Engel's parameter has no absolute scale.

\subsection{The intrinsic pencil (known)}
The following result is printed with status K as prescribed. The intrinsic-pencil and base-preservation assertions for the standard data are \cite[Theorem 7.1]{DGKXv1}; the canonical bundle is \cite[Main Theorem (4); Theorem 9.1]{AC26}. The all-affine-data scope and direct-image formulation are retained from FAM A.1.1; the scope of the cited theorem is distinguished in the source table.
\begin{cited}[FAM A.1.1; K]\label{fam:A.1.1}
For the constructed holomorphic map $f_c:X_c\to\mathbb P^1$, and also for every admissible choice of its affine translations and cusp collar, we have
\[
 \omega_{X_c}\simeq\mathcal O_{X_c}(-2S_2),\qquad
 f_{c*}\mathcal O_{X_c}(kS_2)=\mathcal O_B(\lfloor k/4\rfloor p_2)\quad(k\in\mathbb Z).
\]
The latter equality is an equality of subsheaves of meromorphic functions under pullback. The anticanonical section ring is $\mathbb C[U,V]$, with $\deg U=1$, $\deg V=2$. The complete degree-two map is $f_c$ followed by a projective coordinate isomorphism of its base. Every biholomorphism between two such threefolds is over the identity of $B$.
\end{cited}
\uses{Definition~\ref{def:conventions}; \cite[Theorem 7.1]{DGKXv1}; \cite[Theorem 9.1]{AC26}.}
\begin{cited}[DGKX Theorem 7.2; K]\label{dgkx:all-maps-factor}
For every $c\in\mathcal P$, every holomorphic map
$h:X_c\to\mathbb P^1$ factors uniquely as
$h=\varphi\circ f_c$ with a holomorphic map
$\varphi:\mathbb P^1\to\mathbb P^1$.
\end{cited}
\uses{\cite[Theorem 7.2, \texttt{thm:all-maps-factor}]{DGKXv1}.}

\subsection{The elliptic surface and its reference lift}
\begin{proposition}[FAM B.1.1; T]\label{fam:B.1.1}
Let $E_\tau=\mathbb C/(\mathbb Z+\tau\mathbb Z)$, $P=[\mu]$, and $L_\mu=\mathcal O_{E_\tau}(P-O)$. The holomorphic map
\[
 [\xi,\eta]\longmapsto[\xi,e^{2\pi i\eta}]
\]
identifies $\mathbb C^2/\Pi_c\Lambda$ with the quotient of $L_\mu^\times$ by the lift
$A(\xi,v)=(\xi+6\mu,e^{2\pi i\beta_c}v)$. Engel's generator, in Definition 2.3 of the attached paper, is its inverse, with translation $-6P$ and multiplier $e^{-2\pi i\beta_c}$. The section $P=[\mu]$ is a generator of the Mordell--Weil group of the stated pointed Weierstrass surface.
\end{proposition}
\uses{Definition~\ref{def:conventions}; FAM A.3.1 and preceding dictionary identities; \cite[Remark 3.12]{AC26}; \cite[Lemma 2.1; Corollary 2.2; Definition 2.3]{E26}.}
\begin{proof}
\par\smallskip\noindent\textbf{Step 1.} First quotient by columns two, three and four. Exponentiation kills column four; columns three and two then give
\[
 (\xi,v)\sim(\xi+1,v)\sim(\xi+\tau,e^{2\pi i\mu}v).
\]
For the odd theta function with periods $1,\tau$, the quotient $\theta(\xi-\mu)/\theta(\xi)$ is a meromorphic section of this line bundle: its factors are $1,e^{2\pi i\mu}$ and its divisor is $P-O$. Sending its divisor section to this quotient is the asserted line-bundle isomorphism. Quotienting the remaining first column gives exactly $A$. On sufficiently small simply connected charts the logarithm of $v$ is a holomorphic inverse; changes of logarithm are column-four periods. Thus the displayed map is a biholomorphism of the two quotient complex spaces, with its inverse explicitly defined by local logarithms. Choosing $A^{-1}$ changes the generator of the same infinite cyclic group and gives precisely Engel's declared inverse-translation convention.

\par\smallskip\noindent\textbf{Step 2.} For clarity the elliptic surface in this dictionary is the displayed cubic, rather than only a curve with the same $j$-value. Let $f_{\rm sc}$ be the branch satisfying
\[
 f_{\rm sc}^4=\frac{12t^3(t-1)}{g_2(\tau)},\qquad
 f_{\rm sc}^6=-\frac{8t^4(t-1)^2}{g_3(\tau)},\qquad
 f_{\rm sc}/t\longrightarrow i\sqrt3/\pi.
\]
The identities $g_2^3/(g_2^3-27g_3^2)=t$ make these equations compatible. Their zero and pole orders give a holomorphic branch on the period cover; the indicated cusp sign fixes it. The map $\xi\mapsto(f_{\rm sc}^2\wp_\tau(\xi),\frac12f_{\rm sc}^3\wp'_\tau(\xi))$ is a pointed isomorphism to
$y^2=x^3-3t^3(t-1)x+2t^4(t-1)^2$. Its factors are $f_{\rm sc}(g_1z)=-f_{\rm sc}/\tau$, $f_{\rm sc}(g_2z)=f_{\rm sc}/\tau$; therefore this isomorphism descends over $B^\circ$.

\par\smallskip\noindent\textbf{Step 3.} We first verify that $P=[\mu]$ gives a rational section of the resolved elliptic surface. On $B^\circ$ its cubic coordinates descend by the displayed scale and affine laws. At the two finite centres, $\mu_j$ is not a lattice point, while $f_{\rm sc}$ has orders two and one in the root parameter. Its cubic coordinates are therefore holomorphic and invariant, hence holomorphic in the base coordinate. At the cusp $\mu$ is holomorphic and bounded, and the convergent multiplicative expansion of $\wp$ gives meromorphic cubic coordinates with finite pole orders. Even when its exponential meets one, the pole is controlled by its finite vanishing order; it cannot be identically one because the two finite values of $\mu$ differ. The coordinates are thus meromorphic on $\mathbb P^1$ and rational. The closure of its graph on the proper resolved surface is an irreducible curve; its normalization maps finitely with degree one to the smooth base, so is the base itself. This gives the unique section of that surface extending $P$.

\par\smallskip\noindent\textbf{Step 4.} We also verify that the specified section is a Mordell--Weil generator using its full logarithm. A rational section of this cubic pulls back to a holomorphic section of the good elliptic family at both finite covering points: multiply its projective coordinates by a common power to remove poles, and use the nonsingular limiting Weierstrass equation. At the cusp put $r_{\rm cusp}=1/t$, $X=r_{\rm cusp}^2x,Y=r_{\rm cusp}^3y$. Its equation is $Y^2=X^3-3(1-r_{\rm cusp})X+2(1-r_{\rm cusp})^2$. A section cannot pass through the node $(1,0)$: substitution of $X=1+O(r_{\rm cusp})$, $Y=O(r_{\rm cusp})$ leaves the coefficient $+1$ of $r_{\rm cusp}$ on the left-hand side of the defining equation. At every other central point the multiplicative uniformization is a local isomorphism. Explicitly, at $r_{\rm cusp}=0$ it is
\[
 X=1+\frac{12U}{(U-1)^2},\qquad
 Y=\frac{12\sqrt3 U(U+1)}{(U-1)^3},\qquad U=e^{2\pi i\xi},
\]
with inverse $U=(Y/(X-1)+\sqrt3)/(Y/(X-1)-\sqrt3)$ away from $O$; a local parameter at $O$ is $X/Y$. The convergent lattice expansion of $\wp$ gives the same holomorphic map for small $r_{\rm cusp}$. Hence a section has a holomorphic unit $U(r_{\rm cusp})$ and a bounded logarithm at the cusp. Lifting the section on the simply connected cover gives a global logarithm $\nu$, with this cusp branch. Its only remaining ambiguity is an integer, since a nonzero multiple of $\tau$ is unbounded.

\par\smallskip\noindent\textbf{Step 5.} Its two affine laws have integral defects:
\[
 \nu(g_1z)=-\nu/\tau+\alpha_1+\beta_1\tau(g_1z),\quad
 \nu(g_2z)=\nu/\tau+\alpha_2+\beta_2\tau(g_2z).
\]
Composing them and using cusp invariance gives $\alpha_2=-\beta_1$, $\beta_2=\alpha_1+\beta_1$. Put $n=-\alpha_1-2\beta_1$ and $a=\alpha_1+\beta_1$. Substitution shows that $n\mu+a$ has exactly these defects. Their difference is a global section of the first-coordinate sheaf $\mathcal O_B(-p_0)$ from A.3.1, so is zero. Every rational section is therefore $nP$. The same cusp argument and the distinct fixed values of $\mu$ show that $P$ has infinite order. This proves $\operatorname{MW}(S)=\mathbb ZP$.
\end{proof}
\begin{proposition}[FAM B.1.2; T]\label{fam:B.1.2}
The exact set of nonzero parameters for which every ordinary fibre and both normalized finite root fibres of Engel's period-reference quotient are compact complex two-tori is
\[
 0<|u|<e^{-2\pi M}.
\]
This is the parameter domain for the compact family with both smooth finite root models. The condition on ordinary fibres alone at equality is determined by whether a maximizing point of $D_{\partt}$ is ordinary; this boundary issue is treated in Section~\ref{sec:periods}. For an independent reference, the radius is $|\lambda_0|e^{-2\pi M}$.
\end{proposition}
\uses{Definition~\ref{def:conventions}; FAM A.3.1 and preceding dictionary identities; \cite[Remark 3.12]{AC26}; \cite[Lemma 2.1; Corollary 2.2; Definition 2.3]{E26}.}
\begin{proof}
\par\smallskip\noindent\textbf{Step 1.} The function $N(\xi,v)=|v|\exp(2\pi\operatorname{Im}\mu\,\operatorname{Im}\xi/\operatorname{Im}\tau)$ is invariant under both line-bundle period identifications. The inverse lift multiplies it by $e^{2\pi D_c}$. More directly the real determinant of the four period columns is a nonzero sign times $(\operatorname{Im}\tau)D_c$. Thus the four columns generate a full lattice exactly when $D_c\ne0$, and then the quotient is a compact torus by a closed fundamental parallelepiped. At $D_c=0$ their real span has dimension three; a rank-four lattice in $\mathbb R^4$ cannot result, so the quotient cannot be a compact complex two-torus. The function $D_{\partt}$ descends continuously to the connected base with the cusp removed and tends to $-\infty$ at the cusp. Its image is therefore $(-\infty,M]$: the maximum is attained, and the intermediate value theorem fills all smaller values. Consequently $D_c$ is nonzero at every base point exactly when $\operatorname{Im}c<-M$. Since $|u|=e^{2\pi\operatorname{Im}c}$, the assertion follows, with the indicated reference scaling.
\end{proof}
The geometric identification in this proposition is already implicit in \cite[Remark 3.12]{AC26}. Its intersection calculation re-derives the height computation of \cite[Main Theorem, entry 49; (3.2)]{OS91} and \cite[Theorem 8.6; (8.16)]{Shi90} in the given charts.
\begin{proposition}[FAM B.1.3; T]\label{fam:B.1.3}
For the original affine logarithm $\mu$, we have $\wp_\tau(\mu)=0$ identically. Equivalently the generator $P=[\mu]$ of B.1.1 is one of the two sections $P_*=(0,\sqrt2t^2(t-1))$ and $-P_*$. This assertion does not use a numerical evaluation of $\mu$ or of $\beta$.
\end{proposition}
\uses{Definition~\ref{def:conventions}; FAM A.3.1 and preceding dictionary identities; \cite[Remark 3.12]{AC26}; \cite[Lemma 2.1; Corollary 2.2; Definition 2.3]{E26}.}
\begin{proof}
\par\smallskip\noindent\textbf{Step 1.} We give the intersection calculation that identifies the generator. At $t=1$, put $v_t=t-1$ and blow up the origin. In the $v_t$-chart $x=v_t X,y=v_t Y$, the strict equation is
$Y^2=v_t X^3-3(1+v_t)^3X+2(1+v_t)^4$. Its exceptional curve is smooth, and $P_*$ meets it at $(X,Y)=(0,\sqrt2)$. In the other chart $v_t=xU,y=xY$, the equation is $Y^2=x-3(1+xU)^3U+2(1+xU)^4U^2$. The implicit function theorem solves for $x$ at the common point; the two fibre branches $U=0,x=0$ meet with intersection multiplicity two. These charts cover the tangent conic and resolve the singularity. The two components have multiplicity one; the one met by $P_*$ is the nonidentity component $D$, with $D^2=-2$ from intersection with the full fibre.

\par\smallskip\noindent\textbf{Step 2.} At $t=0$, two blow-ups in the successive $t$-charts give
\[
 Y^2=tX^3+3t^2(1-t)X+2t^2(1-t)^2,
\quad
 Y_2^2=t^2X_2^3+3t(1-t)X_2+2(1-t)^2.
\]
The second equation has two smooth multiplicity-one components $E_{2,\pm}$ at $Y_2=\pm\sqrt2$, and $P_*$ meets one of them. The remaining singular point, in the second $X$-chart, has
\[
 V^2=H(X,T),\quad H=X^2T+(2+3X)T^2-(4X+3X^2)T^3+2X^2T^4,
 \qquad t=XT.
\]
Let $T_0(X)$ solve $H_T(X,T_0)=0$. It has $T_0=-X^2/4+O(X^3)$. Taylor's formula gives
$H-H(X,T_0)=(T-T_0)^2h_0$, $h_0(0,0)=2$, and $H(X,T_0)=X^4e(X)$, $e(0)=-1/8$. These identities follow by substitution in the displayed polynomial, not by discarding its higher terms. With $Z=(T-T_0)\sqrt{h_0}$, $a=V-Z$, $b=(V+Z)/e(X)$, the equation becomes exactly $ab=X^4$ through a local analytic isomorphism. Its resolution has charts
\[
 a=U_i^{i+1}V_i^i,\quad b=U_i^{3-i}V_i^{4-i},\quad X=U_iV_i
 \quad(0\le i\le3),
\]
with changes $U_{i+1}=1/V_i,V_{i+1}=U_iV_i^2$. They give a chain of three $(-2)$-curves. These are the charts of the two successive point blow-ups resolving $ab=X^4$, so the resolution is proper. Retaining $t=XT$ in the comparison gives multiplicities $2,3,2$ on that chain. The two new multiplicity-one components meet its ends; the first exceptional curve meets its middle component, and the original strict fibre meets the other end of that first exceptional curve. All these statements follow by substituting the displayed monomials and their orders; the only fibre components of multiplicity one are the strict component and $E_{2,\pm}$.

\par\smallskip\noindent\textbf{Step 3.} Order the six nonidentity components by the graph
$\Theta_1-\Theta_2-\Theta_3-\Theta_4$ and $\Theta_2-\Theta_5-\Theta_6$, with $E_{2,\pm}=\Theta_4,\Theta_6$. Their intersection matrix is $-C$, where
\[
 C=\begin{pmatrix}2&-1&0&0&0&0\\-1&2&-1&0&-1&0\\0&-1&2&-1&0&0\\0&0&-1&2&0&0\\0&-1&0&0&2&-1\\0&0&0&0&-1&2\end{pmatrix},\quad
 C^{-1}=\frac13\begin{pmatrix}6&9&6&3&6&3\\9&18&12&6&12&6\\6&12&10&5&8&4\\3&6&5&4&4&2\\6&12&8&4&10&5\\3&6&4&2&5&4\end{pmatrix}.
\]
Multiplication verifies this inverse. Self-intersections follow as well from the complete multiplicity vector and $F\cdot\Theta_i=0$.

\par\smallskip\noindent\textbf{Step 4.} The form $dx\wedge dt/y$ is holomorphic and nonvanishing on the finite resolved charts. In the last charts its expression is a unit times $dX\wedge da/a=-dU_i\wedge dV_i$; the earlier blow-up equations cancel the apparent denominators at $y=0$. At infinity it is $-r_{\rm cusp}^{-1}dX\wedge d r_{\rm cusp}/Y$ and has one fibre pole. Thus $K_S=-F$. Adjunction gives $Q^2=-1$ for every section $Q$, since $F\cdot Q=1$. A section meets a unique multiplicity-one component transversely at a smooth fibre point: the pullback of the full base equation to the section has order one, excluding all other possibilities.

\par\smallskip\noindent\textbf{Step 5.} Let $v_Q$ be the intersection vector with $\Theta_i$ and let $d_Q=Q\cdot D$. The possibilities are $v_Q=0,e_4,e_6$ and $d_Q=0,1$. Define the rational divisor class
\[
 \varphi(Q)=Q-O-(Q\cdot O+1)F+\sum\Theta_i(C^{-1}v_Q)_i+\tfrac12d_QD.
\]
It is orthogonal to $O,F,\Theta_i,D$ by direct multiplication. The chord-and-tangent rational function on the generic cubic has divisor $(Q)+(R)-(Q+R)-(O)$; extending it to the surface adds only vertical components. Applying the displayed orthogonal projection therefore proves $\varphi(Q+R)=\varphi(Q)+\varphi(R)$. Expansion gives
\[
 h(Q)=-\varphi(Q)^2=2+2Q\cdot O-v_Q^tC^{-1}v_Q-\tfrac12d_Q.
\]
For $Q\ne O$, the nonnegative intersection $Q\cdot O$ gives $h(Q)\ge2-4/3-1/2=1/6$. The visible $P_*$ is disjoint from $O$ everywhere, and the explicit blow-up charts place it on both nonidentity components. Hence $h(P_*)=1/6$. The logarithm argument in B.1.1 already proved $P_*=nP$ for an integer $n\ne0$, without identifying the visible section. Additivity gives $1/6=n^2h(P)\ge n^2/6$, so $n=\pm1$. The pointed cubic map in B.1.1 has $x(P)=f_{\rm sc}^2\wp_\tau(\mu)$ with $f_{\rm sc}\ne0$ on the ordinary locus. This proves the asserted zero identity.
\end{proof}
\begin{proposition}[FAM B.1.4; T]\label{fam:B.1.4}
The line-bundle comparison of the period family with Engel's family over $B^\circ$ transports its inverse period lift to a global section of $\operatorname{Hom}(t_{6P}^*M_S,M_S)$ with divisor $F_\infty$. Take this transported section for $c=0$ as the reference linearization. Then $\lambda_0=1$ exactly and $u=e^{-2\pi ic}=w^2$. For another independently fixed reference the precise constant is its scalar ratio with this section. With the cusp branch fixed in Section~\ref{sec:periods} and $f_{\rm sc}/t\to i\sqrt3/\pi$, the section is $P=(0,-\sqrt2t^2(t-1))$.
\end{proposition}
\uses{Definition~\ref{def:conventions}; FAM A.3.1 and preceding dictionary identities; \cite[Remark 3.12]{AC26}; \cite[Lemma 2.1; Corollary 2.2; Definition 2.3]{E26}.}
\begin{proof}
\par\smallskip\noindent\textbf{Step 1.} The zero identity $\wp_\tau(\mu)=0$, proved without numerical input in B.1.3, identifies $P$ with the signed visible section. With the cusp branch in Section~\ref{sec:periods} and the specified scale, substitution of $U=-(5+2\sqrt6)$ in the nodal uniformization of B.1.1 gives $Y=-\sqrt2$, so $P=(0,-\sqrt2t^2(t-1))$.

\par\smallskip\noindent\textbf{Step 2.} To compare the line bundles as families, use their common divisor $P-O$ and rigidify at the zero section. Put $F_\theta(\xi)=\theta(\xi-\mu)/\theta(\xi)$, $r_\theta=\theta(-\mu)/\theta'(0)$, and let $s_{P-O}$ be the canonical meromorphic divisor section. The map
\[
 [\xi,v]\longmapsto\frac{2r_\theta}{f_{\rm sc}}
        v\frac{s_{P-O}(\xi)}{F_\theta(\xi)}
\]
is a holomorphic line-bundle isomorphism, including at $P,O$ after cancellation of the equal divisors. At $O$ it sends the chosen frame to $dx/y$, since $F_\theta\sim r_\theta/\xi$ and $x/y\sim-\xi/f_{\rm sc}$ with the corresponding canonical divisor frame. Two such maps that agree at $O$ differ by a holomorphic scalar on each compact elliptic fibre and hence agree everywhere. This proves descent under the actual period identifications.

\par\smallskip\noindent\textbf{Step 3.} Engel's \cite[Lemma 2.1]{E26}, which we have read together with its proof, identifies $\operatorname{Hom}(t_{6P}^*M_S,M_S)$ with $\pi_S^*\mathcal O_B(1)$. The lemma applies because the preceding pointed cubic map and Mordell--Weil calculation identify his $(S,O,P)$ with ours, including the chosen sign. To determine which section is transported, it suffices to compute its order at one generic smooth point of each special fibre, since a section of this pulled-back line is a base section. At a finite centre write $6\mu=p+r\tau+\delta$ with $(p,r)=(4,-2),(3,-3)$ and $\delta/f_{\rm sc}$ holomorphic; this divisibility follows by linearizing the affine laws, whose homogeneous orders are two and one. Evaluate at $\xi=f_{\rm sc}z$ for generic fixed $z$. The factor $(2r_\theta/f_{\rm sc})/F_\theta(f_{\rm sc}z)$ tends to $2z$, a unit. Restoring the lattice shift $p+r\tau$ at the target changes the fibre multiplier to $e^{-2\pi i(\beta_c-r\mu)}$, whose finite limits are $e^{-2\pi i\kappa_1}$ and $e^{-2\pi i\kappa_2}$. The source and target theta ratios remain units. Therefore the base section has order zero at both finite points. At the cusp $f_{\rm sc}$ has a simple pole and $r_\theta$ has a nonzero limit; the two factors $2r_\theta/f_{\rm sc}$ cancel between source and target. The remaining multiplier is
$e^{-2\pi i\beta_c}=q_\tau e^{-2\pi i(\beta_c+\tau)}$, of exact order one. Thus the transported section has divisor $F_\infty$. Engel's \cite[Corollary 2.2]{E26} now identifies it with a nonzero scalar multiple of his reference, since two base sections with the same divisor have constant ratio. Changing $\beta$ by $c$ multiplies it exactly by $e^{-2\pi ic}$. With our stated period-reference normalization, this scalar is one. For any independently chosen reference $\phi_{\rm ref}$, the scalar is instead the precisely defined ratio $\lambda_0=\phi_{\rm part}/\phi_{\rm ref}$, evaluated on any ordinary fibre. Scaling the reference by $a\in\mathbb C^*$ changes $\lambda_0$ to $\lambda_0/a$; no scalar value is invariant under that declared choice.
\end{proof}
\subsection{Correction to Engel's Remark 8.1}
\noindent\textbf{Printed sentence.} In \cite[Remark 8.1]{E26}: ``We may thus build a proper, holomorphic submersion $\mathcal Y^*\to\Delta^*$ whose fiber over $u$ is $X(u,a_0,a_1)$''. The fault is the global descent for fixed sphere logarithmic data; the local construction is retained. No contact with the source's author is recorded. The finite order-two obstruction below identifies the precise qualification.
\begin{proposition}[FAM B.2.3; T]\label{fam:B.2.3}
The assertion in \cite[Remark 8.1]{E26} of a global proper family over the original punctured $u$-disc for fixed sphere logarithmic data requires a double-cover qualification. The retained family with such data does not descend through $u=w^2$. It does give a proper holomorphic family on each simply connected $u$-open set, and a global proper family over the punctured $w$-disc as in A.3.1.
\end{proposition}
\uses{FAM A.2.1, A.3.1, A.4.1 and B.2.1; \cite[Remark 8.1]{E26}.}
\begin{proof}
\par\smallskip\noindent\textbf{Step 1.} Descent through $u=w^2$ would furnish an isomorphism between the two pulled-back fibres at $w$ and $-w$, namely $X_c$ and $X_{c+1}$. A.2.1 excludes that isomorphism by the nonzero class of $d$ in $\Lambda^{A_2}/N_2\Lambda$. This is an exact order-two obstruction on the normalized finite cover. It applies to all sphere data by B.2.1 because $r_2$ is odd.

\par\smallskip\noindent\textbf{Step 2.} There is also a direct analytic witness near zero. The relative Albanese invariant of the multiplicity-four reduced fibre in a holomorphic descended family would be a single-valued holomorphic function of $u$. After pullback it would be even in $w$. But A.4.1 gives $j_2(w)=C^{-1}w^{-1}+744+O(w)$, where $C=e^{-\pi i\kappa_2^{\rm part}}\ne0$; its Laurent coefficient of $w^{-1}$ is nonzero. Thus it is not even on any punctured neighbourhood. This verifies the obstruction independently of a choice of isomorphism lift. Local inverse logarithms of $u$ give the proper local families already constructed; coherent even shifts give their global double-cover family. The untwisted middle quotient itself has the integer shear identification, so this refutation concerns the fixed sphere logarithmic data in the stated remark.
\end{proof}

\section{Discrete data and parameter identifications}\label{sec:discrete}
\subsection{The standard data (known)}
The $2\mathbb Z$ equivalence is Theorem 1.4 of Du, Guo, Kusakabe and Xie, arXiv:2610.00250, version 1 \cite[Theorem 1.4]{DGKXv1}. Version 2 does not contain it. The same statement was obtained independently in the workbench from which this paper grew, and Appendix A prints that proof in the present notation.
\begin{cited}[FAM A.5.1; K]\label{fam:A.5.1}
Fix the original cusp-regular period solution, the clockwise peripheral marking, the prescribed periodic triangular cusp fan, the affine vectors $(\varepsilon,-\varepsilon')$, and the untwisted cusp collar. For every $c,c'\in\mathcal P$, the constructed compact complex threefolds satisfy
\[
 X_c\simeq X_{c'}\quad\Longleftrightarrow\quad c'-c\in2\mathbb Z.
\]
Every biholomorphism is over the identity of the original $\mathbb P^1$. The proper holomorphic family of their isomorphism classes has parameter $0<|w|<e^{-\pi M}$ and nonzero Kodaira--Spencer class at every point.
\end{cited}
\uses{FAM A.1.1, A.2.1, A.3.1 and A.4.2; \cite[Theorems 1.3, 1.4 and 7.1; punctured-disc corollary]{DGKXv1}.}

\subsection{All affine vectors and cusp sections}
\begin{proposition}[FAM B.2.1; T]\label{fam:B.2.1}
Let
\[
 w_1=r_1\varepsilon+k_1d,\qquad w_2=r_2\varepsilon'+k_2d,
 \qquad 3\nmid r_1,\quad 2\nmid r_2,
\]
and use an arbitrary holomorphic cusp collar section $\sigma$ in the construction of \cite[Section 6.4]{AC26}. Write $l_0=\gamma(c(\sigma))$ in its positive cusp-winding convention, and set $p=12l_0-4r_1-3r_2$. Let primes denote a second such datum with the same period solution and cusp fan. Then
\[
 X(c;w_1,w_2,\sigma)\simeq X(c';w'_1,w'_2,\sigma')
\]
holds exactly for the following arithmetic conditions:
\[
 n=c'-c\in\mathbb Z,\qquad
 p'=a p\text{ for some }a\in\{1,-1\},\qquad
 n\equiv k_2-k'_2\pmod2.
\]
All maps are over $B$. A map with the indicated $a$ has fibre-linear part $aI_2$. In particular the complete invariants are
\[
 \bigl(|p|,\ e^{-\pi i(c+k_2)}\bigr),\qquad
 |p|\ge1,\quad\gcd(|p|,6)=1.
\]
Every value of this positive integer occurs. The full collar section enters the construction, but its additional analytic part and its second integral winding contribute no further invariant.
\end{proposition}
\uses{Definition~\ref{def:conventions}; FAM A.1.1--A.3.1 and the preceding discrete-data calculation; \cite[Section 6.4]{AC26}; \cite[Sections 3 and 5]{E26}.}
\begin{proof}
\par\smallskip\noindent\textbf{Step 1.} The determinant calculation in A.1.1 was independent of translations, so the base is still intrinsic. The centralizer proof gives $c'=c+n$ and $aI_2$ exactly as before. Put $S_n=I+nE$ and compare all target periods in source coordinates. On the finite root cover the exact deck equation is
\[
 B_j(g_jz)-R_jB_j(z)=\Pi_c(g_jz)\left(\frac{d_j}{m_j}+\ell_j\right),
 \quad d_j=S_nw'_j-a w_j,
 \quad N_j\ell_j=-d_j.                                      \tag{B2.1}
\]
Necessity follows by the normalized-base-change lift already proved. Conversely the norm equation has a rational solution $r_j^*$ of
$(I-A_j)r_j^*=d_j/m_j+\ell_j$, since $\ker N_j$ over $\mathbb Q$ equals the image of $I-A_j$. We can choose $\gamma(r_j^*)=0$: the columns $\widehat u,\widehat w$ already span that image over $\mathbb Q$. Then $B_j=\Pi_cr_j^*$ is holomorphic on the whole root disc and solves the exact equation. Thus the norm condition is sufficient for a local isomorphism, not only for conjugacy on the central real torus.

\par\smallskip\noindent\textbf{Step 2.} The computed norm images translate (B2.1) into
\[
 r'_1\equiv ar_1\pmod3,\quad r'_2\equiv ar_2\pmod4,
 \quad ak_2-k'_2-nr'_2\equiv0\pmod2.                       \tag{B2.2}
\]
There is no condition on $k_1$. Since $r'_2$ is odd, the last condition is the displayed parity. On the finite collar the middle translation is
$A_j=B_j+(\log s_j)/(2\pi i)\,\Pi_cd_j$ and has integral defect $\ell_j$. Its $\gamma$ defect is
$\alpha_j=(ar_j-r'_j)/m_j$.

\par\smallskip\noindent\textbf{Step 3.} We prove the precise cusp necessity used next. A holomorphic isomorphism between cusp neighbourhoods with fibre-linear part $aI_2$ has, after lifting near a generic smooth central point to the toric covers, the form
\[
 (x,t_c)\longmapsto(t_c^q v_{\rm unit}(t_c)x^a,t_c),\qquad q\in\mathbb Z^2,
 \quad v_{\rm unit}(0)\in(\mathbb C^*)^2.                           \tag{B2.3}
\]
Here the assertion concerns this local lift near a generic central point. To prove it, take a source chart near the generic point of its component and a target chart near its image. The etale toric quotient permits a local lift over a simply connected neighbourhood. The torus characters in a chart at the height-one ray $v$ are $x_i=t_c^{v_i}X_i$, with $X_i$ invertible. On $t_c\ne0$ the affine fibre equation gives $x'_i/x_i^a$ independent of the $X$ coordinates. Its expression in the two charts is therefore $t_c^{v'_i-av_i}$ times a holomorphic unit of $t_c$, including zero. This proves (B2.3). Taking logarithms shows that its winding is in $\Lambda_{\rm tor}=\mathbb Z\widehat w+\mathbb Zd$. Changing a lift adds $(I-A_0)\Lambda=\Lambda_{\rm tor}$, so the vanishing of its class in $\bar\Lambda=\Lambda/\Lambda_{\rm tor}$ is intrinsic. This proof supplies the necessary extension test; no assertion about arbitrary punctured translations is used.

\par\smallskip\noindent\textbf{Step 4.} We retain the actual collar convention. The map from the cusp piece to the middle family is $T_{-\sigma}G_0^{-1}$. Consequently the middle translation induced by (B2.3) is
\[
 A_0=sq+a\sigma-\sigma'+\eta_0,
\]
with $\eta_0$ holomorphic on the cusp disc after a logarithm of the unit. For the clockwise generator $s\mapsto s-1$, the class of its defect in $\bar\Lambda$ is $-a c(\sigma)+c(\sigma')$. On the other hand every middle integral cocycle satisfies
$\ell_0=-A_0(\ell_1+A_1\ell_2)$, whose $\gamma$ coordinate is $-\alpha_1-\alpha_2$. Equality of these coordinates is exactly
\[
 l'_0-a l_0=-(ar_1-r'_1)/3-(ar_2-r'_2)/4,
\]
or $p'=ap$. This proves the remaining necessity. Notice that $p\equiv-r_1\pmod3$ and $p\equiv r_2\pmod2$, so $p$ is nonzero and coprime to six.

\par\smallskip\noindent\textbf{Step 5.} For sufficiency, $p'=ap$ implies the first two congruences of (B2.2) by reduction modulo three and four. The parity supplies the remaining one. Choose integral solutions $\ell_1,\ell_2$ of (B2.1). Their $\gamma$ coordinates already have the required values. Adding any element of $\ker N_1=(I-A_1)\Lambda$ preserves (B2.1). In particular adding an arbitrary integer multiple of $e_1=(0,1,-2,0)^t$ changes the $\widehat u$ coordinate of $\ell_0$ by an arbitrary integer while leaving its $\gamma$ coordinate fixed. We can therefore arrange
$[\ell_0]=-a c(\sigma)+c(\sigma')$ in the entire rank-two group $\bar\Lambda$. Choose lifts of $\sigma,\sigma'$ on the cusp half-plane. Their remaining defect differs from $\ell_0$ by an actual toric lattice vector; an integral $q$ makes $sq+a\sigma-\sigma'$ have exactly defect $\ell_0$.

\par\smallskip\noindent\textbf{Step 6.} The finite translations and this cusp reference now define a locally nonempty affine torsor with linear difference sheaf precisely $\mathcal V$ of A.3.1. The terms $a\sigma-\sigma'$ are retained, even when their chosen lifts have essential singularities at the puncture; they belong to the specified collar reference, while differences of allowed extensions are holomorphic $\eta_0$. The proof of $H^1(B,\mathcal V)=0$ gives a global middle translation. The finite maps descend by (B2.1), and the cusp map is $x\mapsto t_c^qv_{\rm unit}(t_c)x^a$. The fan is preserved by all integer translations and by $y\mapsto-y$; exponentials of the integer period shift agree. Hence this cusp map and its inverse extend over every stratum. The finite overlap is the equality defining $A_j$, and the cusp overlap is the equality defining $A_0$. Thus the maps glue globally.

\par\smallskip\noindent\textbf{Step 7.} Finally parity is equivalent to equality of the two displayed exponentials. Passing from $p$ to $-p$ is effected by $a=-1$ and the same proof; for fixed sign the congruences have just been shown sufficient. Every positive integer coprime to six is obtained by choosing $r_1\in\{1,2\}$ and odd $r_2$ in the appropriate residue classes modulo three and four, and then taking $l_0=(p+4r_1+3r_2)/12$. This also proves that the extra collar data have been accounted for, rather than omitted from the statement.
\end{proof}
\begin{proposition}[FAM B.2.2; T]\label{fam:B.2.2}
For multiplicities $(3,4)$ and $|p|=1$, every admissible discrete member is biholomorphic to one member of A.5.1, uniquely modulo $2\mathbb Z$. At fixed nonzero admissible Engel parameter $u$ there are exactly two isomorphism classes, exchanged by a positive circuit of $u$ about zero. In the convention that Engel's root generator is the inverse of our clockwise one and his circle class is $d$, the identity of the middle quotient sends
\[
 X_{c+1}\longleftrightarrow X(u;a_0-d/3,\ a_1+d/4).
\]
Reversing both declared root generators reverses these two displayed signs.
\end{proposition}
\uses{Definition~\ref{def:conventions}; FAM A.1.1--A.3.1 and the preceding discrete-data calculation; \cite[Section 6.4]{AC26}; \cite[Sections 3 and 5]{E26}.}
\begin{proof}
\par\smallskip\noindent\textbf{Step 1.} Take the canonical target data $(r'_1,r'_2,k'_1,k'_2,l'_0)=(1,-1,0,0,0)$, for which $p'=-1$. Choose $a$ so that $p'=ap$, and an integer $n\equiv k_2\pmod2$. B.2.1 supplies the isomorphism to $X_{c+n}$. The uniqueness is A.5.1. At fixed $u=e^{-2\pi ic}$, its logarithms differ by integers; the invariant $e^{-\pi i(c+k_2)}$ takes exactly the two square roots of $u$. Both occur, already as $X_c,X_{c+1}$, and A.2.1 distinguishes them.

\par\smallskip\noindent\textbf{Step 2.} We spell out the integral comparison to Engel's common basis $B_0=(e_1,e_2,e_3,e_4)$, retaining his $C f_4=e_1+e_4$:
\[
 P_{\rm mark}=\begin{pmatrix}0&-2&-1&-2\\2&-1&0&-1\\0&2&1&3\\-1&0&0&0\end{pmatrix},\qquad
 T_0P_{\rm mark}=P_{\rm mark}A_1,\quad T_1P_{\rm mark}=P_{\rm mark}A_2.
\]
Its determinant is $-1$, $P_{\rm mark}d=\delta_0$, $\psi P_{\rm mark}=-\gamma$, $P_{\rm mark}\varepsilon=-e_4$ and $P_{\rm mark}\varepsilon'=\delta_0-Cf_4$. These are exact lattice homomorphism identities, verified independently through Viaclovsky's intermediate basis in \texttt{ab\_geometry.py}. They also agree with the complex periods: at $z_1$ the complex-linear map
$(\zeta_1,\zeta_2)\mapsto(-\zeta_1/(\rho+1),\zeta_2+\zeta_1/3)$ carries $\Pi_cP_{\rm mark}^{-1}$ to Engel's columns $(1,0),(\rho-1,0),((1+\rho)/3,1/3),(0,-\kappa_1)$. At $z_2$ his $Cf_4$ has preimage $d-\varepsilon'$ and fixed period $1-\kappa_2$. Thus the comparison fixes the positive circle class and the actual local complex structures, in addition to their monodromy matrices.

\par\smallskip\noindent\textbf{Step 3.} For the sign assertion, inverting the concrete clockwise affine generator changes its invariant translation $\Pi_cv_j/m_j$ to $-\Pi_cv_j/m_j$. Thus our dictionary is $a_{j-1}=-P_{\rm mark}v_j/m_j$ in Engel's lattice marking (and its period image in the torus). Integer shift gives $\Pi_{c+1}v_j=\Pi_c(v_j+\gamma(v_j)d)$, with $\gamma(v_1)=1,\gamma(v_2)=-1$, proving the two signs. A positive $u$ circuit changes $\log u$ by $2\pi i$ and hence changes $c$ by $-1$, exchanging the square roots. This gives the claimed permutation with its orientation.
\end{proof}

\subsection{Coprime higher multiplicities}
The theorem concerns the stated model and marking. The identification with the coordinates of \cite[Section 2]{VL26} is separately announced below; this distinction retains the gap noted by the supplied review.
\begin{proposition}[FAM B.3.1; T]\label{fam:B.3.1}
Let $k,l$ be positive integers with $\gcd(3k,4l)=1$, and put $m_1=3k,m_2=4l$. Consider Viaclovsky's sphere members expressed in the same clockwise marking, with $w_j=r_j\varepsilon_j+s_jd$, $\gcd(r_j,m_j)=1$, cusp winding $l_0$, and $p=m_1m_2l_0-m_2r_1-m_1r_2=\pm1$, where $\varepsilon_1=\varepsilon,\varepsilon_2=\varepsilon'$. Two such members are biholomorphic precisely when their multiplicity pair agrees and, for some $a\in\{1,-1\}$ and integer $n=c'-c$,
\[
 p'=ap,\quad r'_j\equiv ar_j\pmod{m_j},\quad
 nr'_1\equiv as_1-s'_1\pmod{k},\quad
 nr'_2\equiv as_2-s'_2\pmod{2l}.
\]
The congruences on $r'_j$ already follow from $p'=ap$. For fixed discrete vectors and cusp collar the identification subgroup is $2kl\mathbb Z$. For fixed admissible $u$ and this multiplicity pair, all the stated sphere translation data give exactly $2kl$ isomorphism classes, permuted in one cycle by a circuit of $u$ about zero.
\end{proposition}
\uses{FAM A.1.2, A.2.1, B.2.1; \cite[Proposition 4.7]{VL26}.}
\begin{proof}
\par\smallskip\noindent\textbf{Step 1.} We first recover the base for these sphere members without using the degree-two canonical pencil. A biholomorphism $F:X\to X'$ gives $g=f'F$. Since $H^2(S^6,\mathbb R)=0$, the closed form $g^*\omega_{\rm FS}$ is exact. Restrict to a regular torus $T$ and wedge with a flat Kahler form on $T$. Stokes' theorem gives integral zero, while the integrand is nonnegative and strictly positive wherever $dg|_T\ne0$. Thus $g$ is constant on each regular fibre. The image of $(f,g)$ is a compact analytic curve whose projection to $B$ has degree one; its normalization gives a holomorphic $h:B\to B$ with $g=hf$. Apply the same construction to $F^{-1}$ to see that $h$ is an automorphism. The three homological local monodromy orders are three, four and infinity, independently of $k,l$; their intertwining maps under $F$ therefore fix each special point. Hence $h=1$. The nonreduced fibre multiplicities now force equal $k,l$. This is the full reduction in the read Proposition 4.7 of Viaclovsky's paper, written here with its exact maps.

\par\smallskip\noindent\textbf{Step 2.} The root model of degree $m_1$ is the pullback of the degree-three model by $s\mapsto s^k$, and likewise at the other point. Its linear action remains $A_1$ or $A_2$, while its root generator has order $m_j$. The primitive invariant coordinate $r_j/m_j$ proves freeness of every nonidentity power. The normalization and deck-commutation arguments of A.1.2 and A.2.1 apply to these root covers without changing a generator. Their norms are
\[
 N_{3k}=kN_1,\qquad N_{4l}=lN_2,
\]
so the fixed-coordinate norm images are $3k\mathbb Z\varepsilon+k\mathbb Zd$ and $4l\mathbb Z\varepsilon'+2l\mathbb Zd$. The norm equation (B2.1), with $m_j$ replaced by these actual covering degrees, gives exactly the displayed congruences. The centralizer and first-column comparison remain unchanged.

\par\smallskip\noindent\textbf{Step 3.} For completeness the integer $p$ records the primitive sphere relation in this marking. The common torus coinvariants are detected by $\gamma$, since $(A_1-I)\Lambda+(A_2-I)\Lambda=\ker\gamma$ by their explicit saturated bases. The finite meridian relations are $x^{m_1}=c^{r_1}$, $y^{m_2}=c^{r_2}$, and the cusp relation is $xy=c^{l_0}$ with $c$ central. Eliminating $y$ gives an abelian relation matrix with determinant $-p$ and primitive first elementary divisor because $\gcd(m_1,m_2)=1$. Thus the sphere data have $p=\pm1$. Reducing $p'=ap$ modulo $m_1,m_2$ gives the two congruences for $r_j$.

\par\smallskip\noindent\textbf{Step 4.} To prove sufficiency, use the local rational solutions of the norm equation and the identical cusp comparison in B.2.1. The first-coordinate extension on the new root discs has vanishing orders $2k$ and $l$, precisely the orders of the pulled-back $\Theta$. Dividing by it again gives $\mathcal O_B(-p_0)$, and finite averaging again gives the same rank-two extension sheaf. Thus $H^1(B,\mathcal V)=0$ still removes the analytic cocycle. The images of $I-A_j$ and their arbitrary $\widehat u$ adjustments are unchanged, so $p'=ap$ is exactly the remaining cusp winding condition. Every local map and overlap is the one already written in B.2.1, with the new denominator $m_j$; this proves global sufficiency.

\par\smallskip\noindent\textbf{Step 5.} For identical data, $p\ne0$ excludes $a=-1$. The congruences reduce to $k\mid nr_1$ and $2l\mid nr_2$. The two primitivity assumptions imply $k\mid n$ and $2l\mid n$. Since $\gcd(k,2l)=1$, their least positive common multiple is $2kl$. For different sphere translation data, $r'_1$ is invertible modulo $k$, $r'_2$ modulo $2l$, and the Chinese remainder theorem gives exactly one integer shift class modulo $2kl$. The values of $s_1$ modulo $k$ and $s_2$ modulo $2l$ all occur without affecting freeness or $p$. Thus there are exactly $2kl$ classes at fixed $u$. The change $c\mapsto c-1$ on a positive $u$ circuit acts by a generator of this cyclic residue group. Program \texttt{ab\_exact.py} verifies the repeated norms and shift steps for five coprime pairs; the preceding symbolic identities prove all pairs.
\end{proof}
\begin{announcement}[PUSH E.1.1; A]\label{push:E.1.1}
Let $(\delta,a,b,d)$ be the marking in
\cite[Section 2, equations (2.1), (2.2), (2.7), (2.11)--(2.12)]{VL26},
and put $u_0=d+2a+2b$, $u_1=d+3a+3b$.
The homomorphism from the marking of the question to this marking is
\[
 P_V=\begin{pmatrix}
 -2&1&0&1\\-6&0&-1&0\\0&-1&0&0\\-1&0&0&0
 \end{pmatrix},\qquad \det P_V=-1.
\]
It intertwines the two finite monodromies and satisfies
\[
 P_V\widehat\delta=\delta,\qquad
 P_V\varepsilon=-u_0,\qquad
 P_V\varepsilon'=\delta-u_1.
\]
With the counterclockwise deck generators of that paper, its data
$m_0=3k$, $m_1=4l$ and
$a_i=[v_i/m_i]$ with $v_i=\rho_i\delta+N_i u_i$ correspond to
\[
 w_1=N_0\varepsilon-\rho_0\widehat\delta,
 \qquad
 w_2=N_1\varepsilon'-(\rho_1+N_1)\widehat\delta.
\]
For the section-preserving cusp comparison used there, $l_0=0$ and
$p=-(nN_0+mN_1)$.  Thus its sphere condition is
$b:=nN_0+mN_1=\pm1$.
The brackets in $a_i$ mean the torsion point in the smooth
root torus, using its section-preserving origin and the
transported period vector.  Replacing $v_i$ by an integral
multiple $m_i$ times a period does not change that torsion
point, although its chosen lift fixes the collar bookkeeping.
\end{announcement}
\uses{FAM B.3.1; PUSH E.1.1--E.1.2; \cite[Section 2]{VL26}.}
\noindent\textit{Mechanism.} The integral change of basis transports the positive-root translations; repeated norm images give the displayed congruences. Proof in the repository, S6\_02\_PUSH, E.1.1 \cite[run S6\_02\_PUSH, E.1.1]{Repo}.
\begin{announcement}[PUSH E.1.2; A]\label{push:E.1.2}
Use a local logarithm $c$ of the transported linearization,
$u=e^{-2\pi i c}$, and the same reference section and cusp comparison
on both sides.  Two sphere members with data
$(m,\rho_0,N_0;n,\rho_1,N_1;c)$ and primed data are biholomorphic
if and only if their multiplicities agree and there are
$\alpha\in\{1,-1\}$ and $h=c'-c\in\mathbb Z$ such that
\[
 b'=\alpha b,\qquad
 hN'_0\equiv\rho'_0-\alpha\rho_0\pmod{k},\qquad
 hN'_1\equiv\rho'_1-\alpha\rho_1\pmod{2l}.
\]
Every such biholomorphism covers the identity of $B$.
For fixed discrete data the period identification subgroup is
$2kl\mathbb Z$.  At fixed $u$ all sphere translation data give
exactly $2kl$ classes.  A positive $u$-circuit acts on them by a
single cycle.
\end{announcement}
\uses{FAM B.3.1; PUSH E.1.1--E.1.2; \cite[Section 2]{VL26}.}
\noindent\textit{Mechanism.} The integral change of basis transports the positive-root translations; repeated norm images give the displayed congruences. Proof in the repository, S6\_02\_PUSH, E.1.2 \cite[run S6\_02\_PUSH, E.1.2]{Repo}.
